\chapter*{Introduction}
\addcontentsline{toc}{chapter}{Introduction}
\setcounter{chapter}{0}

In recent years, automatic fall detection has become an important research topic in Ambient Assisted Living because falls may have severe consequences, particularly for older adults. Interest has consequently grown in monitoring systems that can be used within healthcare and care
facilities and in domestic environments.

The problem nevertheless remains complex, because some daily activities, such
as sitting down or bending over, may appear visually similar to a fall.
For this reason it is not sufficient to consider only the person's final
posture, but also the way in which that posture is reached. It therefore
becomes important to describe the temporal evolution of the body, taking into
account information such as changes in orientation, velocity, relative
displacements of the different anatomical parts and other quantities related
to the dynamics of the movement.

The objective of this thesis is to analyse and verify the possibility of
detecting a fall and distinguishing it from a normal daily activity by means
of a system based on \emph{pose estimation} with the aim of enabling earlier detection, supporting timely assistance and reducing the need to transmit identifiable visual information.

Most existing studies report classifier performance on controlled datasets. However, these results do not necessarily indicate how the complete system behaves when the subject, viewpoint, acquisition domain and computational platform change. This thesis therefore investigates both the relative role of movement representation and classifier architecture, and the progressive
degradation or preservation of performance when the system is moved towards realistic operating conditions.

In more detail, the thesis is divided as follows. Chapter \ref{cap:contesto}
provides an introduction to the field of automatic fall detection, focusing on
the importance of the problem in the older population and on the consequences
a fall entails, on the main requirements a system must satisfy in order to be
genuinely installable, and on the difficulties that arise when it has to
operate in a real environment.
Chapter \ref{cap:stato_arte} reviews the main concepts of fall detection
proposed in the literature, such as the different technologies employed, the
representations of human movement, the models used for temporal
classification and the problems related to datasets and evaluation protocols.
Chapter \ref{cap:metodologia} describes the methodology adopted, following the
path that leads from pose estimation to the final decision. It presents the
geometric representations considered, the construction of the temporal
features, their analysis and selection, and the formation of the sequences
used by the classifiers.
Chapter \ref{cap:protocollo} defines the protocol used for the evaluation and
presents the main experimental results. After class imbalance, subject-wise
partitioning, normalisation and the choice of the decision threshold, the
chapter compares the different representations and architectures and studies
the effect of the viewpoint, the definition of the task and the transfer to an
external dataset.
Finally, Chapter \ref{cap:edge} analyses the implementation of the entire
pipeline on an \emph{edge} device and its validation under conditions closer
to those of real use. It examines in depth the computational cost of the
individual stages, the behaviour of the system on the device and the results
obtained during the sessions recorded in domestic environments, moving from
the individual predictions of the classifier to the alarm events produced
during continuous operation.
Chapter 
ef{cap:discussion} discusses the results obtained as a whole, relating the individual experiments to one another in order to identify which factors most affect the behaviour of the system, and draws the conclusions of the work together with the limitations of the evaluation.
